\scanpage{36}
\mspage{10}
avec les pôles $0,\infty$. D'autre part, l'antipodisme $\sigma'z=-b/\bar z$ \deleted{\unsure{transpose}}\inserted{intervertit} $0$ et $\infty$, et \deleted{\unsure{point $z\in C$}}\inserted{transforme $C$ en lui} $=$ car $z\in C$ i.e. $z\bar z=b$ implique $\sigma'z\in C$ i.e.
\[
\underbrace{\underbrace{(\sigma'z)}_{(-b/\bar z)}\underbrace{\overline{(\sigma'z)}}_{(-b/z)}}_{b^2/z\bar z}=b
\]
Ainsi, $\sigma'$ et $(\sigma,\{0,\infty\})$ définissent la $=$ structure riemannienne sphérique. On est content!

2) $N$ est ext. \uline{non centrale} de $\gamma$ par $T$.

Cela signifie que si $\sigma\in N\setminus T$, on a
\[
\sigma.z=a\bar z,\quad\deleted{\unsure{a}}\quad(a\in\mathbb C^*).
\]
\deleted{\unsure{Donc}} on conclut
\[
\sigma^2z=a(\bar a z)=a\bar a z
\]
et la condition $\sigma^2\in T$ signifie $a\bar a=1$ d'où $a\in\mathbb U$. qui implique aussi \deleted{$\sigma^2=1$}
\[
\sigma^2=1.\qquad\text{i.e. l'extension }N\text{ est \uline{scindée}}
\]
Le produit $u_\theta\sigma$ $(\theta\in\mathbb U)$ est
\[
u_\theta\sigma(z)=(\theta a)\bar z
\]
de la $=$ forme.

Soit $\alpha$\inserted{$\in\mathbb U$} tel que
\[
a=\alpha^{\unsure{2|\sim}}=\alpha/\bar\alpha\qquad\text{donc}
\]
\[
z\in\mathbb C^\sigma\Longleftrightarrow z=\alpha\frac{\bar z}{\unsure{\bar z|\bar\alpha}}\Longleftrightarrow z/\alpha=\overline{(z/\alpha)}\quad\text{i.e.}\quad z\in\alpha\mathbb R
\]

\margin{Ici $N$ est déterminé en termes de $T$, i.e. en termes de $\{z,z'\}$, comme $T\cup N^-$, où $N^-$ est formé de \uline{toutes} les anti-involutions de $X$ qui fixent $\{z,z'\}$. Donc il n'est pas question, dans cette sorte de \unsure{données}, d'extraire une métrique riem. sphérique.}
% archive foliation: 38
